% AUTO-GENERATED by the update_record tool. DO NOT EDIT.
% verify_paper regenerates this file from record.json and the run outputs
% and fails on any difference, so manual edits always surface.
\noindent\textbf{Hypothesis 1.} LLM-AutoSciLab の仮説条件付き獲得（候補機構の予測が最も食い違う実験を選ぶ）は、著者らが評価したスカラー応答のオラクル（NewtonBench / ActiveSciBench）の外でも機構回復を改善する。具体的には、全種の濃度時系列が観測され SBML を提出する SciGym-small の反応回復課題で、同じ gpt-6-luna を主モデルに使ったとき、SciGym 公式の ReAct エージェント（実験を自由記述の推論で選ぶ）より低い STE と高い RMS を出し、その改善は実験選択則によるものである。~\cite[s2.p2, s2.p3, s2.p5, s2.p8, s2.p14]{kabra-2026-llm}, \cite[s1.p3, s1.p8]{duan-2025-measuring}
\begin{enumerate}
\item[\textbf{Claim 1}] SciGym-small 137 件の平均 STE（trajectory\_smape）は、LLM-AutoSciLab の SciGym 用ループ（autoscilab-luna）の方が SciGym 公式 ReAct エージェント（react-luna）より 0.02 以上低い。~\cite[s2.p1, s2.p2, s2.p6, s2.p9, s2.p10, s2.p11, s2.p12, s2.p13, s2.p15, s2.p16, s2.p19, s2.p22, s2.p23, s2.p24]{kabra-2026-llm}, \cite[s1.p4, s1.p5, s1.p10, s1.p12]{duan-2025-measuring}
  \emph{Rationale:} h1 の中心。同じ主モデル・同じ実験予算・同じ 137 件・同じ公式 Evaluator で、手法だけを替えたときの差が STE に出れば、仮説条件付きの閉ループが SBML の機構回復にも効くという h1 の主張の直接の証拠になる。
  \emph{Criterion:} \texttt{\detokenize{autoscilab-luna}}.\texttt{\detokenize{trajectory_smape}} $-$ \texttt{\detokenize{react-luna}}.\texttt{\detokenize{trajectory_smape}} $\leq -0.02$.
  \emph{Prediction:} $[-0.1, -0.02]$ (同じ openai/gpt-6-luna の公式 ReAct は auto-res2/scigym-frontier-models（137 件、max\_iterations=20）で STE 0.2295 だった。論文は同じ予算で 2–5 倍のサンプル効率（s2.p1）を報告しているが、SBML への翻訳とフィットはこちらのアダプタなので、改善幅は 0.02〜0.10 と控えめに見積もる。).
  \emph{Observed:} 0.098266.
  \emph{Verdict:} refuted.
\item[\textbf{Claim 2}] 同じループで実験選択だけを仮説と無関係にした autoscilab-nohyp-luna の平均 STE は、autoscilab-luna より 0.01 以上高い。~\cite[s2.p7, s2.p11, s2.p22]{kabra-2026-llm}
  \emph{Rationale:} h1 の「改善は実験選択則によるもの」の部分。論文で最も効いた部品（s2.p7）を外した run との差が出れば、c1 の差が仮説生成やフィットではなく仮説条件付き獲得に帰せられる。
  \emph{Criterion:} \texttt{\detokenize{autoscilab-nohyp-luna}}.\texttt{\detokenize{trajectory_smape}} $-$ \texttt{\detokenize{autoscilab-luna}}.\texttt{\detokenize{trajectory_smape}} $\geq 0.01$.
  \emph{Prediction:} $[0.01, 0.08]$ (論文の Figure 4 では仮説条件付き獲得を外すと全ベンチマークで大きく落ちる（s2.p7）。SciGym では実験 20 回のうち情報の薄い点が増える分の差として 0.01〜0.08 と見積もる。).
  \emph{Observed:} 0.0047988.
  \emph{Verdict:} refuted.
\item[\textbf{Claim 3}] SciGym 公式 ReAct エージェント（react-luna）の RMS（修飾子なしの reaction\_f1）は、autoscilab-luna より 0.02 以上低い。~\cite[s1.p1, s1.p2, s1.p3, s1.p4, s1.p6, s1.p10, s1.p11]{duan-2025-measuring}
  \emph{Rationale:} h1 の RMS の部分。STE は軌道の近さしか見ないので、反応の構造が実際に回復されているかは RMS で確かめる。react-luna はこの claim の run であると同時に c1 の比較対象でもある。
  \emph{Criterion:} \texttt{\detokenize{react-luna}}.\texttt{\detokenize{reaction_f1}} $-$ \texttt{\detokenize{autoscilab-luna}}.\texttt{\detokenize{reaction_f1}} $\leq -0.02$.
  \emph{Prediction:} $[-0.15, -0.02]$ (同じ openai/gpt-6-luna の公式 ReAct は auto-res2/scigym-frontier-models で reaction\_f1 0.3594 だった。SciGym 論文 Table 1 の RMS F1（修飾子なし）はモデル間で 0.10〜0.34 の幅（s1.p10）なので、手法の差として 0.02〜0.15 を見積もる。).
  \emph{Observed:} -0.0609217.
  \emph{Verdict:} supported.
\end{enumerate}
\noindent\emph{Assumed, so that the claims together imply Hypothesis 1:}
\begin{itemize}
\item STE / RMS / NTS（SciGym 公式 Evaluator、airas-eval scigym\_small）が機構回復の良さを表す（c1, c2, c3）
\item アダプタが SciGym の設定を LLM-AutoSciLab の手法に写す仕方（仮説 = 不完全モデルに足す反応の集合 + 速度則、実験 = 初期濃度の変更、食い違い = 候補 SBML の濃度時系列の対数分散）が手法の意図を保っている（c1, c2, c3）
\item 上流コード（scientific-discovery/LLM-AutoSciLab @351aae7、h4duan/SciGym @88a7b93）はそれぞれの論文を実装している（c1, c2, c3）
\item knob の一覧（design の arguments）は、コードから列挙できた範囲を超えて完全である（c1, c2, c3）
\item Vercel AI Gateway が返すモデル（openai/gpt-6-luna、meta/llama-3.1-8b）はその名前のモデルである（c1, c2, c3）
\item LLM のサンプリングは非決定的であり、1 回の run の差は再現時にずれうる（c1, c2, c3）
\item GitHub Actions の実行記録と git 履歴は正しい（c1, c2, c3）
\item ReAct の max\_iterations=20 と AutoSciLab のオラクル問い合わせ 20 回を同じ実験予算とみなす（ReAct の反復にはコード実行だけの反復も含まれるので、ReAct の実験回数は 20 以下になる）（c1, c3）
\end{itemize}
